\documentclass{article}
\usepackage{amsmath,amssymb,amsthm,mathtools}
\usepackage{leanify}
\newtheorem{theorem}{Theorem}
\newtheorem{lemma}{Lemma}
\title{Reasons, restatements and chosen objects in proofs}
% A test of proof sentences that give their reasons in words, restate an informal sentence,
% or choose an object: `INFORMAL, that is, FORMAL`; `since WORDS` at the start and in the
% middle of a sentence, with a relative clause; a description before `so`; `Let $c>0$ be
% such that …`; `we can find … such that …`; `we obtain a NP $X$ such that …`; headings of
% paragraphs; a section as a citation; `for each of the WORDS $x\in S$`; `Letting
% $\varepsilon\to0$ gives the claim.` The expected output of `leanify check --claims` is
% check.txt.
\begin{document}
\maketitle

\section{Squares}\label{sec:squares}

Every square of a real number is nonnegative, that is, for every real number $x$, we have $x^2\ge0$.
We have $x^3=x$ for each of the three numbers $x\in\{-1,0,1\}$.

\begin{lemma}\label{lem:sum}
For all real numbers $x$ and $y$, we have $x^2+y^2\ge2xy$.
\end{lemma}

\begin{proof}
\emph{Expanding the square.}
We have $(x-y)^2=x^2-2xy+y^2$.
Hence $(x-y)^2\ge0$ by Section~\ref{sec:squares}, and since the square is the expansion above, $x^2+y^2-2xy\ge0$.
\end{proof}

\begin{lemma}\label{lem:archimedes}
For every real number $\varepsilon>0$, there is a positive integer $n$ such that $\frac1n<\varepsilon$.
\end{lemma}

\begin{proof}
By the Archimedean property of the real numbers, we can find a positive integer $n$ such that $n>\frac1\varepsilon$.
Then $\frac1n<\varepsilon$.
\end{proof}

\begin{theorem}\label{thm:zero}
Let $a$ be a real number.
Assume that $a\le\varepsilon$ for every $\varepsilon>0$.
Then $a\le0$.
\end{theorem}

\begin{proof}
Let $\varepsilon>0$.
Since the hypothesis applies to $\varepsilon$, which is positive, $a\le\varepsilon$.
Letting $\varepsilon\to0$ gives the claim.
\end{proof}

\begin{lemma}\label{lem:between}
For every real number $x$, there is an integer $k$ such that $k\le x<k+1$.
\end{lemma}

\begin{proof}
Taking the integer part of $x$, we obtain an integer $k$ such that $k\le x<k+1$.
Then $k\le x$.
\end{proof}

\begin{lemma}\label{lem:small}
Let $b$ be a real number with $0<b<1$.
Then $b^2<b$.
\end{lemma}

\begin{proof}
Let $c>0$ be such that $b+c=1$.
The number $c$ is the distance from $b$ to $1$, so $b^2=b(1-c)$.
Since $c$ is positive, $b^2<b$.
\end{proof}

\end{document}
